%% P09 --- Determinant, inverse and change of basis
%%
%% Every number in this program that the reader cannot do in their head comes
%% from code/p09_determinant_inverse.py and is pulled in with \val{}. The
%% console listing is a committed transcript written by that script.
%%
%% THE THESIS: the determinant is a signed volume scale factor; zero means
%% information was destroyed, which is Program P08's null space seen from the
%% other side; and an inverse exists exactly when nothing was.
%%
%% WHAT THIS PROGRAM IS OWED, read out of the files rather than remembered:
%%   P04  ALREADY DOES CHANGE OF BASIS AS A WORKED EXAMPLE -- the point
%%        (3, 4) against axes turned by 30 degrees, with the coordinates
%%        changing and the length not. So the IDEA is spent. What is left is
%%        better: the change is itself a MATRIX, and it is invertible exactly
%%        because nothing was destroyed.
%%   P06  gives composition and non-commutativity, so det(AB) = det(A)det(B)
%%        is a statement about composing scale factors rather than an identity
%%        to check.
%%   P08  hands this program over by name: "rank says how many directions
%%        survive; the determinant says how much of the space does, in one
%%        number that is zero exactly when something was destroyed -- which is
%%        this program's null space seen from the other side." So the det = 0
%%        frame is a GATE against P08's rank rather than a new fact.
%%   F08  has rotation matrices, that a rotation preserves length, and rotary
%%        position embedding in full. So this program adds only the reading
%%        F08 could not give: a rotation is an invertible CHANGE OF BASIS.
%%
%% WHAT THIS PROGRAM LEAVES ALONE, checked against tools/programs.json:
%%   eigenvalues, the spectral theorem, positive definiteness        -> P10
%%   the SVD, the condition number, the pseudoinverse                -> P11
%%     -- so HOW CLOSE to singular is P11's question and not this one, and
%%     the manifest's "measured in P10" was stale by one: P10 undertakes no
%%     such measurement and P11 undertakes "the concrete form of do not
%%     invert". The operation count IS measured here; the numerical reason
%%     is P11's.
%%
%% Twin: programs/pl/P09-determinant-inverse.tex. Frame for frame, macro for
%% macro, number for number; tools/parity.py compares them.

\program{Determinant, inverse and change of basis}
\label{prog:P09}
\index{determinant}
\index{matrix!inverse}

\begin{outcomes}
\outcome{1--7}{say what one number can measure about a whole map, and why the
  sign of it is not decoration}
\outcome{8--14}{say what a determinant of zero means about the inputs, and
  connect it to what Program~\ref{prog:P08} counted}
\outcome{15--19}{say what happens to determinants along a product, and why one
  bad factor settles the whole chain}
\outcome{20--27}{say when a map can be undone, and price undoing it against
  answering the question you actually had}
\outcome{28--37}{read a change of basis as a matrix, and say what an orthogonal
  one keeps}
\end{outcomes}

\begin{quiz}
\item A matrix doubles every $x$ and triples every $y$. What happens to the
  area of a region it is applied to? \teachesat{1--2}
  \answerto{It is multiplied by $\val{p09.stretch.det}$. That factor is the
  same for every region, which is what makes one number enough.}
\item What does a determinant of zero tell you about the inputs?
  \teachesat{9--10}
  \answerto{Two different inputs are now sharing an output, so nothing can
  send the result back. It is Program~\ref{prog:P08}'s null space, measured as
  volume instead of counted as directions.}
\item $A$ is $100 \times 100$ with determinant $1$. Every entry is multiplied
  by $\num{0.1}$. What is the determinant now? \teachesat{13--14}
  \answerto{$10^{-100}$. Each of the hundred rows contributes its own factor
  of $\num{0.1}$, and nothing about the map got worse.}
\item You need $x$ with $Ax = b$, for one $b$. Is it cheaper to solve, or to
  form $A^{-1}$ and multiply? \teachesat{24--25}
  \answerto{Solving, by a factor of about $\val{p09.cost.ratio}$ at
  $n = \val{p09.cost.n}$ \dash{} and the two give exactly the same answer.}
\item What object turns one pair of coordinates for a point into the pair
  measured against different axes? \teachesat{28--29}
  \answerto{A matrix, and an invertible one. The point does not move; only its
  description does, so you must be able to get the old description back.}
\item A matrix has determinant $1$. Must it be a rotation? \teachesat{35--36}
  \answerto{No. A shear has determinant $1$ and stretches almost every vector
  it touches. Preserving area and preserving length are two conditions, and
  neither implies the other.}
\end{quiz}

% ============================================================
\section{What one number can say about a whole map}
\label{sec:P09-signed-area}
\index{determinant!as a signed area}

\begin{fr}
Program~\ref{prog:P08} counted. It asked \emph{how many} directions come out
the other side of a matrix, called the answer the rank, and left one question
standing: \emph{how much}.

Here is the move that answers it. A matrix does not only move points, it moves
whole regions, because a region is made of points. So pick a region, apply the
matrix, and compare the two sizes.

Pick the unit square \dash{} corners $(0,0)$, $(1,0)$, $(1,1)$ and $(0,1)$
\dash{} and the matrix that doubles every $x$ and triples every $y$. What shape
do the four corners land on, and what is its area?
\dotline
\nextframe
\end{fr}

\begin{fr}
\begin{ansblock}
A rectangle $2$ wide and $3$ tall. Area $\val{p09.stretch.det}$.
\end{ansblock}

Now generalise it, because the answer was not about that matrix.

The corner $(1,0)$ is sent to the first column of the matrix and $(0,1)$ to the
second \dash{} that is what multiplying by a matrix does to those two vectors,
and Program~\ref{prog:P06} said so. So the image of the unit square is always
the parallelogram spanned by the two columns.

That area has a name. It is the \textbf{determinant}, written $\det A$ or
$\lvert A \rvert$, and for the matrix above it is $\val{p09.stretch.det}$.
\end{fr}

\begin{fr}
One square is a thin basis for a claim about a whole map, so here is why it is
enough.

Any region can be covered by small squares as finely as you like. The matrix is
linear, so it treats every one of those small squares exactly as it treated the
unit one \dash{} same shape of image, same ratio of areas. Add them back up and
the whole region has scaled by the same factor.

So the determinant is not the area of one particular parallelogram. It is
\textbf{the factor by which this matrix scales every area at once}, and that is
why one number can describe a map of the entire plane.
\end{fr}

\begin{rigourbox}
Covering a region with ever finer squares and adding is exactly the
construction Program~\ref{prog:F13} used to turn a sum into an integral, and
making it rigorous is the same work: it needs the limit to exist and the
covering not to matter. Both are true here and neither is proved in this book.
\end{rigourbox}

\mermaidfig{p09-one-number}{Why one comparison of areas describes a map of the
whole plane.}{the determinant as one scale factor}

\begin{fr}
Now a matrix that does something an area alone cannot see.

Take the one that swaps the two coordinates: it sends $(1,0)$ to $(0,1)$ and
$(0,1)$ to $(1,0)$. The unit square goes to the unit square, so the area is
unchanged.

Something has happened to the plane all the same. What?
\dotline
\nextframe
\end{fr}

\begin{fr}
\begin{ansblock}
It has been turned over. Left and right have swapped.
\end{ansblock}

The determinant records that, and it is why the word is \emph{signed} area:
$\det = \val{p09.flip.det}$ for the swap, not $+1$.

The machinery is the shoelace formula, which computes the area of a polygon
from its corners and gives a \emph{signed} answer \dash{} positive if the
corners are listed anticlockwise, negative if clockwise. Take the unit square,
apply a matrix, run the shoelace formula on the image, and you get the
determinant exactly. On $\val{p09.area.trials}$ random matrices, exactly, with
no rounding anywhere, because the arithmetic is done over fractions.
\end{fr}

\begin{trapbox}
\textbf{\enquote{The determinant is the area, so take the absolute value.}}

That throws away the half of it that matters. $\lvert \det A \rvert$ is the
area; $\det A$ is the area \emph{and} the orientation, and the second is a fact
about the map that no amount of measuring will recover.

It has a consequence you can check: no product of rotations is ever a mirror,
because every rotation has determinant $+1$ and a product of numbers equal to
$+1$ is $+1$. A negative determinant is a different kind of map, not a smaller
one.
\end{trapbox}

\begin{fr}
For a $2 \times 2$ matrix the formula is
\[ \det \begin{pmatrix} a & b \\ c & d \end{pmatrix} = ad - bc \]
In three dimensions the same quantity is a signed \emph{volume}, and in $n$ it
is a signed $n$-volume. The formula grows unpleasant and the meaning does not
change at all, which is why this program never writes it down beyond
$2 \times 2$: a machine will compute it by elimination, and so will you if you
ever have to.

Apply $ad - bc$ to the swap of the two coordinates, which has $a = d = 0$ and
$b = c = 1$.
\dotline
\nextframe
\end{fr}

\begin{fr}
\begin{ansblock}
$0 \times 0 - 1 \times 1 = \val{p09.flip.det}$.
\end{ansblock}

Which is what the shoelace formula said two frames ago, arrived at from the
algebra instead of from the picture. The stretch is worth doing the same way:
$a = 2$, $d = 3$, $b = c = 0$, and $ad - bc$ gives $\val{p09.stretch.det}$.

It is worth being explicit now about what has and has not been claimed, because
the determinant attracts more folklore than any other number in linear algebra.

What has been claimed is one thing: \textbf{the determinant is the factor by
which signed volume scales}. Everything this program does from here is a
consequence of that sentence, including the two that people memorise
separately \dash{} what zero means, and why determinants multiply.

If you find yourself reaching for a rule about determinants that you cannot
derive from that sentence, the rule is probably about something else.
\end{fr}

% ============================================================
\section{Zero, and what it destroys}
\label{sec:P09-zero}
\index{determinant!zero}

\begin{fr}
Take a matrix whose second column is twice its first \dash{} say columns
$(1, 2)$ and $(2, 4)$.

Where does the unit square go?
\dotline
\nextframe
\end{fr}

\begin{fr}
\begin{ansblock}
Onto a line segment. Both columns point the same way, so the parallelogram they
span is flat.
\end{ansblock}

A flat parallelogram has no area, so by the previous section the determinant is
$0$ \dash{} and $ad - bc = 1 \times 4 - 2 \times 2$ agrees.

That is the whole of what a zero determinant is. Not a special case, not a
degeneracy the formula happens to produce: the square came out with nothing
inside it.
\end{fr}

\begin{fr}
And now the consequence, which is the reason this program sits after
Program~\ref{prog:P08} rather than before it.

If the whole plane has been pressed onto a line, then points that were apart
are now together \dash{} two different inputs share an output. No rule can send
that output back, because it would have to return two answers to one argument.

Program~\ref{prog:P08} counted the inputs this happens to and called them the
null space, and said a matrix is invertible exactly when that space holds
nothing but zero. Here it is again as volume: $\det A = 0$ exactly when
$\operatorname{rank} A < n$. Checked on $\val{p09.sing.trials}$ matrices, over
fractions, agreeing every time \dash{} which is not a discovery but a gate,
because the two programs must not be able to disagree about what singular
means.
\end{fr}

\begin{fr}
Here are both numbers from one matrix, computed by the same elimination
Program~\ref{prog:P04} wrote and Program~\ref{prog:P08} reused.

\transcript{p09-singular}

The rows are $1, 2, 3$ then $4, 5, 6$ then $7, 8, 9$, which is about as
unremarkable as a matrix gets. Why is its determinant zero?
\dotline
\nextframe
\end{fr}

\begin{fr}
\begin{ansblock}
The third row is twice the second minus the first. $2 \times (4,5,6) - (1,2,3)
= (7,8,9)$.
\end{ansblock}

So the three rows do not reach out into three dimensions. They lie in a plane,
and a matrix whose rows lie in a plane cannot fill a volume: the unit cube comes
out flat, and the determinant says so with a zero.

The rank of $\val{p09.worked.rank}$ beside it is the same statement counted
rather than measured, which is what the previous frame promised.
\end{fr}

\mermaidfig{p09-zero-means-flat}{A determinant of zero, and the loss it is
reporting.}{what a zero determinant destroys}

\begin{fr}
Now a question whose obvious answer is wrong, and it is worth writing your
answer down before reading on.

$A$ is $100 \times 100$ and $\det A = 1$. Somebody multiplies every entry of it
by $\num{0.1}$.

What is the determinant of the new matrix?
\dotline
\nextframe
\end{fr}

\begin{fr}
\begin{ansblock}
$10^{-100}$.
\end{ansblock}

Each row has been scaled by $\num{0.1}$, and scaling a row scales the volume,
so a hundred rows contribute a hundred factors.

Nothing about the map got worse. It sends independent vectors to independent
vectors exactly as before; every direction that survived still survives; the
rank is unchanged at $100$. A machine solving with it will behave exactly as
well. The only thing that changed is the units.
\end{fr}

\begin{trapbox}
\textbf{\enquote{A small determinant means the matrix is nearly singular.}}

The determinant scales like the $n$th power of the matrix, so in any dimension
worth caring about its size is dominated by scaling and says almost nothing
about how close to flat the map is. $10^{-100}$ above, for a matrix that is not
close to singular in any sense.

The determinant answers \emph{is anything destroyed} exactly, in yes and no.
The question \emph{how nearly is something destroyed} is a different question
with a different answer \dash{} it is the condition number, and it belongs to
Program~\ref{prog:P11}. Reading the first as an approximate answer to the
second is how a working engineer ends up chasing a number that was never about
what they thought.
\end{trapbox}

% ============================================================
\section{Scale factors multiply}
\label{sec:P09-product}
\index{determinant!of a product}

\begin{fr}
$B$ scales every area by $\det B$. Then $A$ scales every area by $\det A$.

Apply $B$ and then $A$ \dash{} which is the matrix $AB$, by
Program~\ref{prog:P06}. By what factor does that scale an area?
\dotline
\nextframe
\end{fr}

\begin{fr}
\begin{ansblock}
$\det A \times \det B$. A scale factor applied to a scale factor is a product.
\end{ansblock}

So $\det(AB) = \det(A)\det(B)$, checked exactly on two hundred random pairs.

It is worth noticing what this did not require. There is no calculation here at
all: the identity is a restatement of what a determinant \emph{is}, and the
only reason it is usually presented as something to verify is that the
$2 \times 2$ formula makes it look like an accident of algebra.

And one thing falls out of it that surprises people. Program~\ref{prog:P06}
made a good deal of $AB \ne BA$ \dash{} but $\det(AB) = \det(BA)$ always,
because both are $\det A \det B$ and numbers commute. The order matters to the
map and not to the volume.
\end{fr}

\begin{fr}
Now put it together with Program~\ref{prog:P08}'s bottleneck.

A chain of matrices is applied one after another, and one of them somewhere in
the middle is singular. Four more full-rank matrices are applied after it.

What is the determinant of the whole chain?
\dotline
\nextframe
\end{fr}

\begin{fr}
\begin{ansblock}
Zero. One factor of zero in a product is the end of it.
\end{ansblock}

Which is Program~\ref{prog:P08}'s result told in volumes rather than in
directions: what a narrow layer discards, nothing downstream can give back. The
script builds a $3 \times 3$ matrix through a $2$-wide middle, confirms it is
singular, and then stacks four full-rank matrices on top of it, checking after
each that the determinant is still exactly zero.

Two programs, two arguments, one fact. Program~\ref{prog:P08} got there by
counting directions and this one by multiplying scale factors, and the second
route is shorter because a product of numbers is easier to reason about than a
composition of maps.
\end{fr}

\begin{fr}
One practical consequence of determinants multiplying, and it is
Program~\ref{prog:P02}'s subject arriving in a new place.

A model that composes many layers has a determinant that is a product of many
determinants. Each is some ordinary number; the product of a few hundred of
them is not, and will underflow to zero or overflow to infinity long before it
means anything.

So nobody computes it. What is carried instead is $\ln \lvert \det \rvert$,
where the product becomes a sum and the sum stays in range \dash{} the same
move Program~\ref{prog:F03} made for the probability of a sequence, for the
same reason, and Program~\ref{prog:P02} says why it is not optional.
\end{fr}

\begin{aibox}
The place you will meet that term by name is a \emph{normalising flow}, whose
whole design is a chain of invertible layers: because each layer can be undone,
a density can be pushed through it, and the change-of-variables formula charges
you $\ln \lvert \det \rvert$ of each layer's Jacobian for the privilege.

That is why flow architectures are built out of layers whose determinant is
cheap \dash{} triangular ones, where it is the product of the diagonal, so the
term costs $n$ multiplications instead of the $n^{3}$ that a general
determinant would.

This book does not teach flows. The term is named here so that when you meet
\code{logdet} in a loss function you know what it is charging for: it is the
volume, and it is a log because the product was not survivable.
\end{aibox}

% ============================================================
\section{The inverse, and why the code does not form it}
\label{sec:P09-inverse}
\index{matrix!inverse}

\begin{fr}
$A$ is square. When is there a matrix $A^{-1}$ that undoes it \dash{} that
sends $Ax$ back to $x$ for every $x$?

Note that the question is only being put to a square matrix, and
Program~\ref{prog:P08} says why it has to be. A wide matrix must send some
non-zero vector to zero, so there is no undoing it; a tall one cannot reach
most of the space it maps into, so for most outputs there is nothing to undo.
Only a square matrix is a candidate at all.

You have every part of the answer already. Say it before reading on.
\dotline
\nextframe
\end{fr}

\begin{fr}
\begin{ansblock}
Exactly when nothing was destroyed: $\det A \ne 0$.
\end{ansblock}

If $\det A = 0$ then two inputs share an output, and an inverse would have to
return both of them. If $\det A \ne 0$ then no two inputs share an output, and
the columns reach the whole space, so every output has exactly one input to be
sent back to.

Four sentences, one condition:

\begin{itemize}
\item $A$ is invertible;
\item $\det A \ne 0$;
\item $\operatorname{rank} A = n$;
\item the null space contains nothing but the zero vector.
\end{itemize}

Program~\ref{prog:P08} wrote the third and fourth and this program adds the
second, which is the one you can check with a single number.
\end{fr}

\begin{fr}
Two small facts, one of which is worth deriving rather than remembering.

The first: $A^{-1}A = I$ exactly, which the script checks over fractions on
sixty matrices, because \enquote{exactly} is a claim and floating point would
not have supported it.

The second: $A$ scales volume by $\det A$ and $A^{-1}$ undoes it. So what is
$\det(A^{-1})$?
\dotline
\nextframe
\end{fr}

\begin{fr}
\begin{ansblock}
$1/\det A$.
\end{ansblock}

From the previous section and nothing else:
\[ \det(A^{-1})\det(A) = \det(A^{-1}A) = \det I = 1 \]

Notice that this is the second time in two sections that a fact about inverses
came out of the product rule in one line. That is the habit worth taking from
here \dash{} most of what is true about determinants is a consequence of their
multiplying, and deriving it takes less time than looking it up.
\end{fr}

\begin{fr}
Now the part that changes what you write.

You have $A$ and one right-hand side $b$, and you want the $x$ with $Ax = b$.
There are two routes. Eliminate down the columns of $A$ and read the answer back
up. Or form $A^{-1}$, then multiply it by $b$.

Both give the same $x$ \dash{} the script checks that they agree exactly over
fractions, because a difference between them would make the comparison about
something else. At $n = \val{p09.cost.n}$, which is cheaper, and by roughly how
much?
\dotline
\nextframe
\end{fr}

\begin{fr}
\begin{ansblock}
Solving, by a factor of about $\val{p09.cost.ratio}$.
\end{ansblock}

$\val{p09.cost.solve}$ multiplications and divisions against
$\val{p09.cost.invert}$, for the same answer to the same question.

Those two numbers are not quoted from a textbook's operation count. Both
routines count their own arithmetic as they do it \dash{} every multiply and
every divide increments a counter \dash{} so the ratio on this page is a
measurement of the code beside it.
\end{fr}

\mermaidfig{p09-solve-or-invert}{Two routes to the same vector, and the work
each one does.}{solving against inverting}

\begin{fr}
The reason is not subtle once you look at what each route builds.

Elimination builds the answer and nothing else: $n$ numbers. Forming the
inverse builds $n^{2}$ numbers, of which you then use $n$ \dash{} the other
$n^{2} - n$ are the answers to right-hand sides nobody asked about.

So the case for the inverse is when you genuinely have many right-hand sides
and they arrive at different times. That case exists and it is rarer than the
code suggests, because a factorisation of $A$ answers repeated right-hand sides
too, and more cheaply.
\end{fr}

\begin{fr}
And there is a second reason, which is stronger than the first and is not
this program's to give.

Forming the inverse is worse \emph{numerically}, not merely slower: the answer
it produces has larger error than the answer elimination produces, from the
same matrix and the same right-hand side. Saying why needs a way to measure how
much a matrix amplifies error, which this program does not have and
Program~\ref{prog:P11} builds.

So the honest form of the advice is two claims with two owners. The operation
count is measured here. The accuracy is Program~\ref{prog:P11}'s, and it is the
half that decides the argument.
\end{fr}

\begin{aibox}
The line to recognise is the closed form of linear regression, written out as
Program~\ref{prog:P08} derived it:
\begin{center}
\code{x = inv(A.T @ A) @ A.T @ b}
\end{center}
against a solver called on the same system. Both compute the least-squares fit.
The first forms an inverse it does not need, and forms it for $A\T A$, which is
the one matrix in the problem you least want to build \dash{} again for reasons
that are Program~\ref{prog:P11}'s.

It appears in a great deal of working code because it is the formula as
written on the page, and the formula on the page is a description of the
answer rather than an instruction for computing it. Those two things come apart
more often than notation admits.
\end{aibox}

% ============================================================
\section{Change of basis}
\label{sec:P09-basis}
\index{basis!change of}

\begin{fr}
Program~\ref{prog:P04} took the point
$(\val{p04.basis.x}, \val{p04.basis.y})$, wrote it against axes turned by
$\val{p04.basis.angle}$ degrees, and got
$(\val{p04.basis.c1}, \val{p04.basis.c2})$ \dash{} coordinates changed, length
unchanged at $\val{p04.basis.len}$.

What it did not say is what performed the change. What kind of object takes one
pair of coordinates for a point and returns the pair measured against different
axes?
\dotline
\nextframe
\end{fr}

\begin{fr}
\begin{ansblock}
A matrix.
\end{ansblock}

Each new coordinate is how far the point goes along one new axis, which is a
dot product with that axis; two axes give two dot products; two dot products
arranged in rows are a matrix times a vector. For axes turned by $\theta$ it is
\[ \begin{pmatrix} \cos\theta & \sin\theta \\
                   -\sin\theta & \cos\theta \end{pmatrix} \]
and the script applies exactly that to Program~\ref{prog:P04}'s point at
Program~\ref{prog:P04}'s angle and requires Program~\ref{prog:P04}'s own two
committed coordinates back. The two programs are quoting one computation, and
if either moves the build says so.
\end{fr}

\begin{fr}
Now the requirement, which is where this section joins the rest of the program.

Nothing moved. The point is where it was; only the description changed. So the
change has to be reversible \dash{} you must be able to recover the old
coordinates from the new ones, or the new ones are not a description of the
same point at all.

By the previous section that puts a condition on one number. Which number, and
what must be true of it?
\dotline
\nextframe
\end{fr}

\begin{fr}
\begin{ansblock}
The determinant, and it must not be zero.
\end{ansblock}

Which is exactly Program~\ref{prog:P04}'s requirement that a basis be
independent, arriving from a different direction: axes that were not
independent would collapse the plane, and a collapsed description cannot be
undone.

Of all the changes of basis available, one family is used far more than the
rest: the rotations.

Turning the axes and leaving the point alone does not change any distance in
the picture, so a rotation cannot be scaling area.

What, then, is its determinant?
\dotline
\nextframe
\end{fr}

\begin{fr}
\begin{ansblock}
$\val{p09.rot.det}$.
\end{ansblock}

Worked exactly, with the $3$--$4$--$5$ triangle so every entry is a fraction and
you can check it by hand:
\[ R = \begin{pmatrix} \tfrac{3}{5} & -\tfrac{4}{5} \\[2pt]
                       \tfrac{4}{5} & \tfrac{3}{5} \end{pmatrix},
   \qquad \det R = \tfrac{9}{25} + \tfrac{16}{25} = \val{p09.rot.det} \]

Two more properties, both checked exactly rather than to a tolerance. Its
transpose is its inverse: $R\T R = I$. And it preserves length \dash{} the
vector $(\val{p04.basis.x}, \val{p04.basis.y})$ has squared length
$\val{p09.rot.len2}$ before and $\val{p09.rot.len2}$ after.
\end{fr}

\begin{fr}
The middle one of those three is the definition, and it deserves its name.

A matrix whose transpose is its inverse is called \textbf{orthogonal}. It is a
compact condition and it buys a great deal: undoing the change costs nothing
but writing the matrix down the other way, and the lengths and angles that
Program~\ref{prog:P05} built are all preserved, because they are built out of
dot products and $R\T R = I$ is precisely the statement that dot products
survive.

That is the sense in which an orthogonal change of basis is free. It is not
cheap, it is \emph{lossless}: no information about the vectors is left behind
in the old coordinate system.
\end{fr}

\begin{fr}
Which lets Program~\ref{prog:F08}'s last section be said properly.

Rotary position embedding gives each position an angle and rotates the query
and the key by it, and Program~\ref{prog:F08} showed the resulting attention
score depends only on the difference of the two positions.

What could not be said there is what the rotation costs the model, because
saying it needs a determinant. A rotation is an orthogonal change of basis: it
is invertible, its determinant is $\val{p09.rot.det}$, and it destroys nothing.
The position information is added by \emph{re-describing} the query and the key
rather than by altering them, so no direction of the embedding is spent on
carrying it.
\end{fr}

\begin{fr}
Last question of the program, and it catches most people.

$\det = 1$ for a rotation, and a rotation preserves length. Here is a matrix
with determinant $1$:
\[ S = \begin{pmatrix} 1 & 1 \\ 0 & 1 \end{pmatrix} \]
Does it preserve length?
\dotline
\nextframe
\end{fr}

\begin{fr}
\begin{ansblock}
No. It sends $(0,1)$ to $(1,1)$, whose length is $\sqrt{2}$.
\end{ansblock}

$S$ is a shear: it slides the plane sideways by an amount proportional to
height. Areas survive, because a parallelogram pushed over keeps its base and
its height. Lengths do not.

And the same collision runs the other way. Reflecting in the $x$-axis has
determinant $\val{p09.reflect.det}$ and preserves every length there is.
\end{fr}

\begin{trapbox}
\textbf{\enquote{Determinant $1$ means it is a rotation.}}

Two independent conditions, and each has a counterexample against the other on
this page: the shear preserves area and not length, the reflection preserves
length and not orientation.

What makes a rotation is $R\T R = I$ \emph{and} $\det R = 1$ \dash{} the first
for lengths, the second for orientation \dash{} and dropping either one gives a
different family of maps. The first alone is the orthogonal matrices, which
include the reflections.
\end{trapbox}

\begin{fr}
The determinant answered the question Program~\ref{prog:P08} left standing.
Rank counts how many directions survive a matrix; the determinant measures how
much of the space does, in one signed number, and that number is zero exactly
when something was destroyed.

Everything else in the program came out of that. Determinants multiply because
scale factors do. An inverse exists exactly when the determinant is non-zero,
because that is what it means for nothing to have been lost. A change of basis
is invertible because the point did not move. An orthogonal one costs nothing
because it does not even change a length.

Program~\ref{prog:P10} asks the next question, and it is a sharper one. Not
\emph{how much} of the space survives, but \emph{which directions} come through
unturned \dash{} only stretched. There are usually a few of them, the amount
each is stretched by is a number, and the determinant turns out to be the
product of those numbers, which is this program's one number taken apart into
its pieces.
\end{fr}

% ============================================================
\begin{summarybox}
\sumitem{1--3}{\result{\textbf{A matrix scales every area by one factor}, because a
  region is a covering of small squares and each is treated alike. That factor
  is the \textbf{determinant}, and the image of the unit square is the
  parallelogram spanned by the columns \dash{} area $\val{p09.stretch.det}$ for
  the matrix that doubles $x$ and triples $y$}.}
\sumitem{4--5}{\result{\textbf{The area is signed}, and the sign is orientation: the
  swap of the two coordinates has determinant $\val{p09.flip.det}$, having
  turned the plane over. Checked against the shoelace formula on
  $\val{p09.area.trials}$ matrices exactly. $\lvert \det \rvert$ is the area;
  $\det$ is the area and the mirror}.}
\sumitem{6--7}{\result{$ad - bc$ in two dimensions, a signed volume in three, a signed
  $n$-volume in $n$. \textbf{Every later claim in the program is a consequence
  of \enquote{the factor by which signed volume scales}}, so a rule you cannot
  derive from that sentence is probably about something else}.}
\sumitem{8--10}{\result{\textbf{Zero means flat}, and flat means two inputs now share
  an output. $\det A = 0$ exactly when $\operatorname{rank} A < n$, checked on
  $\val{p09.sing.trials}$ matrices \dash{} which is Program~\ref{prog:P08}'s
  null space measured as volume instead of counted as directions}.}
\sumitem{11--12}{\result{The worked case: rows $1,2,3$ then $4,5,6$ then $7,8,9$ has
  determinant $0$ and rank $\val{p09.worked.rank}$, because the third row is
  twice the second minus the first. \textbf{Rows lying in a plane cannot fill a
  volume.}}}
\sumitem{13--14}{\result{\textbf{A small determinant does not mean nearly singular.}
  Scaling a $100 \times 100$ matrix by $\num{0.1}$ takes its determinant to
  $10^{-100}$ and changes nothing about the map. \emph{Is anything destroyed}
  is the determinant's question; \emph{how nearly} is the condition number's,
  and it is Program~\ref{prog:P11}'s}.}
\sumitem{15--16}{\result{\textbf{$\det(AB) = \det(A)\det(B)$}, because a scale factor
  applied to a scale factor is a product \dash{} a restatement rather than a
  calculation. And $\det(AB) = \det(BA)$ always, even though
  Program~\ref{prog:P06} showed $AB \ne BA$: order matters to the map and not
  to the volume}.}
\sumitem{17--18}{\result{\textbf{One singular factor makes the whole chain singular},
  which is Program~\ref{prog:P08}'s bottleneck told in volumes: four full-rank
  matrices stacked on a flattened one leave it flat}.}
\sumitem{19}{\result{A product of hundreds of determinants underflows, so what is
  carried is $\ln \lvert \det \rvert$ \dash{} Program~\ref{prog:F03}'s move for
  Program~\ref{prog:P02}'s reason. That term is a normalising flow's
  change-of-variables charge, and it is why such layers are built triangular}.}
\sumitem{20--21}{\result{\textbf{Invertible, $\det \ne 0$, full rank, and a null space
  of nothing but zero are one condition in four sentences}, and the second is
  the one a single number settles}.}
\sumitem{22--23}{\result{$A^{-1}A = I$ exactly over the rationals, and
  $\det(A^{-1}) = 1/\det A$ in one line from the product rule. \textbf{Most of
  what is true about determinants follows from their multiplying}, and deriving
  it is quicker than looking it up}.}
\sumitem{24--26}{\result{\textbf{Solving beats inverting} \dash{}
  $\val{p09.cost.solve}$ operations against $\val{p09.cost.invert}$ at
  $n = \val{p09.cost.n}$, a factor of $\val{p09.cost.ratio}$, counted by the
  routines themselves. Elimination builds $n$ numbers; the inverse builds
  $n^{2}$ of which you use $n$}.}
\sumitem{27}{\result{And the stronger reason is accuracy rather than speed, which is
  Program~\ref{prog:P11}'s to give. \textbf{The formula on the page is a
  description of the answer, not an instruction for computing it} \dash{} which
  is what \code{inv(A.T @ A) @ A.T @ b} is}.}
\sumitem{28--29}{\result{\textbf{A change of basis is a matrix}, gated against
  Program~\ref{prog:P04}'s own committed coordinates: each new coordinate is a
  dot product with a new axis, and two of them arranged in rows are a matrix
  times a vector}.}
\sumitem{30--31}{\result{It must be invertible, because the point never moved and the
  old description has to be recoverable \dash{} which is
  Program~\ref{prog:P04}'s independence requirement arriving as
  \textbf{a determinant that is not zero}}.}
\sumitem{32--33}{\result{\textbf{A rotation has determinant $\val{p09.rot.det}$ and its
  transpose is its inverse}, which is what \emph{orthogonal} means. Lengths and
  angles survive because dot products do, so an orthogonal change of basis is
  \textbf{lossless} rather than merely cheap}.}
\sumitem{34}{\result{Which is what Program~\ref{prog:F08}'s position encoding costs the
  model: nothing. \textbf{Rotating the query and the key re-describes them
  rather than altering them}, so no direction of the embedding is spent on
  carrying a position}.}
\sumitem{35--36}{\result{\textbf{Determinant $1$ and preserving length are two
  conditions.} The shear that adds a point's height to its width has
  determinant $1$ and stretches; the reflection has determinant
  $\val{p09.reflect.det}$ and does not. A rotation is both at once}.}
\end{summarybox}

\canyou

\begin{testexercises}
\item A matrix has first row $(3, 1)$ and second row $(6, 2)$. Compute its
  determinant, and say without further computation whether it has an inverse.
  \answerto{$3 \times 2 - 1 \times 6 = 0$. No inverse: the second column is
  twice the first, so the plane is pressed onto a line and two inputs share an
  output.}
\item $A$ is $4 \times 4$ with $\det A = 3$, and $B$ is $4 \times 4$ with
  $\det B = -2$. Give $\det(AB)$, $\det(BA)$ and $\det(A^{-1})$.
  \answerto{$-6$, $-6$ and $\frac{1}{3}$. The first two are equal because both
  are $\det A \det B$ and numbers commute, even though $AB$ and $BA$ are
  different matrices.}
\item Every entry of a $10 \times 10$ matrix with determinant $2$ is doubled.
  Give the new determinant, and say whether the matrix has become closer to
  singular.
  \answerto{$2 \times 2^{10} = 2048$. It has not become anything: scaling all
  ten rows scales the volume ten times over, and the map sends exactly the same
  directions to exactly the same places, twice as far along.}
\item You have one $A$ and a thousand right-hand sides, arriving one at a time.
  Argue both sides of whether to form $A^{-1}$ once.
  \answerto{For: the inverse is formed once and each solve is then one
  matrix--vector product. Against: a factorisation of $A$ is formed once too,
  each solve is then two triangular substitutions of the same order of cost,
  and the answers are more accurate. The inverse wins on neither count, which
  is why the case for it is narrower than it looks.}
\item Show that the product of two orthogonal matrices is orthogonal.
  \answerto{$(PQ)\T(PQ) = Q\T P\T P Q = Q\T I Q = Q\T Q = I$. So changing basis
  twice by rotations is one change of basis by a rotation, and nothing is lost
  either time.}
\item A matrix sends $(1,0)$ to $(2,1)$ and $(0,1)$ to $(4,2)$. Give its
  determinant, and say what the unit square becomes.
  \answerto{The columns are $(2,1)$ and $(4,2)$, so $\det = 2 \times 2 - 4
  \times 1 = 0$. The square becomes a segment along the direction $(2,1)$: the
  second column is twice the first, so both edges land on one line.}
\end{testexercises}

\begin{furtherproblems}
\item A matrix has two identical rows. Give its determinant without computing
  anything, and say why.
  \answerto{Zero. Two identical rows are not independent, so the rows lie in a
  space of fewer dimensions than the matrix is wide, and the image of the unit
  cube is flat.}
\item $\det(A + B) = \det A + \det B$ is false. Find a counterexample among
  $2 \times 2$ matrices with entries $0$ and $1$.
  \answerto{Let $A$ carry a single $1$ in its top left corner and $B$ a
  single $1$ in its bottom right, every other entry zero. Both have
  determinant $0$, and $A + B = I$ has determinant $1$. Determinants multiply
  along a product; nothing says they add along a sum, and volumes have no
  reason to.}
\item A rotation by $\theta$ has determinant $\cos^{2}\theta + \sin^{2}\theta$.
  Which program's identity is that, and what does it let you conclude for every
  angle at once?
  \answerto{Program~\ref{prog:F08}'s. It is $1$ for every $\theta$, so no
  rotation ever changes an area or an orientation \dash{} which is one
  computation covering the whole family rather than a check at each angle.}
\item Program~\ref{prog:P08} showed a wide matrix must send some non-zero
  vector to zero. What does that say about the determinant of a square matrix
  built by stacking a wide one on a tall one?
  \answerto{Nothing on its own \dash{} the product of a tall and a wide matrix
  can be square and non-singular. But if the shared middle dimension is smaller
  than the outer one, the determinant is zero, because the rank cannot exceed
  the middle and a full-rank square matrix needs all $n$.}
\item Why is the determinant of a triangular matrix the product of its
  diagonal, and why does that make triangular layers attractive in a
  normalising flow?
  \answerto{Elimination has nothing to eliminate, so the pivots are already the
  diagonal entries and their product is the determinant. It costs $n$
  multiplications instead of the general $n^{3}$, so the $\ln \lvert \det
  \rvert$ term in the loss is nearly free.}
\end{furtherproblems}
